% Fifteen Hands: rulebook. Build with lualatex (needs texlive-luatex).
\documentclass[10pt]{article}
\usepackage[paperwidth=5.83in,paperheight=8.27in,
            textwidth=4.5in,textheight=6.6in,hcentering,vcentering]{geometry}
\usepackage[english]{babel}
\usepackage{fontspec}
\usepackage{booktabs}
\usepackage{xcolor}
\usepackage[protrusion=true,expansion=true]{microtype}

\setmainfont{TeX Gyre Pagella}
\newfontfamily\headfont{Ubuntu}[Ligatures=TeX]

\definecolor{ink}{HTML}{1A1A1A}
\definecolor{rule}{HTML}{8A8578}
\definecolor{accent}{HTML}{7A2A2E}

\linespread{1.06}
\setlength{\parindent}{0pt}
\setlength{\parskip}{0.55\baselineskip}
\frenchspacing
\pagestyle{plain}
\color{ink}

% ---------- headings ----------
\newcommand{\sect}[2]{%
  \par\addvspace{1.3\baselineskip}%
  {\headfont\bfseries\large\color{accent}#1\hspace{0.6em}#2\par}%
  {\color{rule}\hrule height 0.4pt}%
  \addvspace{0.5\baselineskip}}

\newcommand{\subsect}[1]{%
  \par\addvspace{0.8\baselineskip}%
  {\headfont\bfseries\normalsize #1\par}%
  \addvspace{0.2\baselineskip}}

% labelled step; label box sizes itself so word labels do not collide
\newlength{\stepw}
\newcommand{\step}[2]{%
  \par\addvspace{0.35\baselineskip}%
  \settowidth{\stepw}{\headfont\bfseries #1\hspace{0.7em}}%
  \ifdim\stepw<1.6em \setlength{\stepw}{1.6em}\fi
  \noindent\makebox[\stepw][l]{\headfont\bfseries\color{accent}#1}%
  \begin{minipage}[t]{\dimexpr\textwidth-\stepw\relax}#2\end{minipage}\par}

% plain bullet for unlabelled lists
\newcommand{\bul}[1]{%
  \par\addvspace{0.35\baselineskip}%
  \noindent\makebox[1.6em][l]{\color{accent}\textbullet}%
  \begin{minipage}[t]{\dimexpr\textwidth-1.6em\relax}#1\end{minipage}\par}

\newcommand{\keyword}[1]{{\headfont\bfseries #1}}

% boxed note
\newcommand{\note}[2]{%
  \par\addvspace{0.7\baselineskip}%
  \noindent\fcolorbox{rule}{white}{%
    \begin{minipage}{\dimexpr\textwidth-2\fboxsep-2\fboxrule\relax}
    \small{\headfont\bfseries #1}\par\vspace{0.25\baselineskip}#2
    \end{minipage}}%
  \par\addvspace{0.7\baselineskip}}

\begin{document}

% ================= title =================
\thispagestyle{empty}
\vspace*{6\baselineskip}
\begin{center}
  {\headfont\bfseries\fontsize{30}{34}\selectfont LIGHT COIN\par}
  \vspace{1.2\baselineskip}
  {\itshape A game of silver, short weight,\\ and a currency everybody needs and nobody trusts\par}
  \vspace{2.5\baselineskip}
  {\headfont\small 5 PLAYERS \quad\textbullet\quad 12 ROUNDS \quad\textbullet\quad 75 MINUTES\par}
  \vfill
  {\small\itshape No single light coin can be told from a good one.\\
   No single light coin makes anybody poorer.\\
   The coinage fails all the same.\par}
  \vspace{3\baselineskip}
\end{center}
\clearpage
\setcounter{page}{1}

% ================= premise =================
\sect{1}{The Coin}

Your city mints its own money, and each of the five great Houses holds the right to strike it. A coin is worth what it weighs in silver, and it passes from hand to hand all day long because everybody trusts that it does.

You can strike a \keyword{light coin}. A little less silver, the same face stamped on it, the difference too small for anybody to feel in their palm. Every House can do this, and the profit is immediate.

One light coin harms nobody. Nor does a hundred. But coin circulates, and light coin circulates faster, and there comes a point at which the traders begin weighing rather than counting, and then refusing altogether. When that happens the money is worthless in every purse in the city, including yours.

The difficulty is that a purse of coin tells you how much silver is missing but never which House it came from, nor whether your own full-weight striking on any particular day was the thing that kept the money good.

\note{The shape of the game}{You win by accumulating Profit, and striking light is the fastest way to accumulate it. But if Trust reaches zero before the twelfth round, \emph{every player loses, including whoever was winning}. The game is an argument about how much bad silver the city can absorb, conducted by people who cannot check each other's work.}

% ================= components =================
\sect{2}{Components}

\begin{tabular}{@{}p{0.30\textwidth}p{0.62\textwidth}@{}}
\toprule
20 white cubes & Truth. Four per player each round. \\
20 black cubes & Light coin. Four per player each round. \\
1 Melt bag & Opaque. Claims go in here. \\
1 Silver bag & Refilled to 16 white and 4 black cubes each round. \\
5 cups & Opaque. Conceals your retained striking cube. \\
5 jars & Opaque, lidded. Your private record. Never opened until the end. \\
5 screens & Or simply play below the table edge. \\
Trust track & 0 to 40, with tiers marked at 31, 21 and 11. \\
Profit & About 200, in 1s and 5s. \\
20 Discredit tokens & Red. Kept in public view. \\
1 Standard card & With five membership slots. \\
10 Warning tokens & Two per player. \\
5 player aids & The challenge table, reproduced at the back. \\
1 First Player marker & Passes clockwise each round. \\
\bottomrule
\end{tabular}

% ================= setup =================
\sect{3}{Setup}

\step{1.}{Set the Trust marker to \keyword{36}.}
\step{2.}{Give each player one white and one black cube, one cup, one jar, one screen and one player aid.}
\step{3.}{Put 16 white and 4 black cubes in the Silver bag. Leave the Melt bag empty.}
\step{4.}{Each player begins with 0 Profit and no Discredits. The Standard card starts empty.}
\step{5.}{Give the First Player marker to whoever most recently signed something without reading it.}

% ================= the round =================
\sect{4}{The Round}

Twelve rounds, five phases each. The first two matter more than they look, because they are what stops anybody working out who did what.

\subsect{Phase 1. The Striking}

Behind your screen, take one white and one black cube in a closed fist. The Melt bag passes in seat order. Reach in, release \emph{exactly one} cube, and withdraw your fist still closed.

\begin{quote}
Release \keyword{black} to strike light. Release \keyword{white} to strike full weight.
\end{quote}

Then, in full view, place the cube still in your hand under your cup without showing it. That retained cube is the exact opposite of the one you played, and it is now the only proof of what you did.

\note{Do this before you touch the Silver bag}{Both cubes must be committed, one in the melt and one under the cup, before anybody draws silver. A player still holding a cube when the silver comes round can swap the two, and a House that struck light has every reason to gamble on it.}

\subsect{Phase 2. The Silver}

The Silver bag holds sixteen white cubes and four black. It passes in seat order.

Reach in, take \emph{one cube without looking at it}, and drop it straight into the Melt bag with your hand closed.

You will never know what you contributed. Neither will anybody else. Roughly one round in five your silver came in adulterated, and there is no arrangement under which you could have refused it, declared it, or been believed about it afterwards.

\note{Why you are not allowed to look}{If you knew what you had drawn, you would have every reason to hold the black back and put a white cube in instead. Bad silver erodes Trust, and Trust is your money. Every House would reason the same way, no adulterated metal would ever reach the melt, and with it would go the uncertainty that makes an assay a judgement rather than an arithmetic exercise.}

\subsect{Phase 3. The Count}

Tip out the Melt bag. Count the \keyword{black} cubes aloud. That number is the round's \keyword{Light Count}.

It counts both deliberate short weight and bad silver nobody ordered, and there is no way to tell how the total divides between them.

\subsect{Phase 4. The Assay}

Beginning with the First Player and going clockwise, the first player who wishes to may call \keyword{one assay}, and only one assay is called per round in total.

\begin{quote}
You may only call an assay if you released a \keyword{white} striking cube this round.
\end{quote}

Name a House and demand their coin on the scales. They lift their cup.

\step{Light}{White under the cup means they played black. They lose this round's Profit and take a \keyword{Discredit}. The challenger gains \keyword{3 Profit}.}
\step{Sound}{Black under the cup means they struck full weight. \emph{The challenger} takes a Discredit.}

A House taking its \keyword{third Discredit} is ruined and out of the game.

\note{Why only the honest may accuse}{You cannot demand that others put their coin on the scales in a round in which your own was short. Policing therefore has a price: to be able to call an assay at all, you must give up the round's larger payout.}

\subsect{Phase 5. Erosion and Reset}

Reduce Trust by the Light Count.

Then, without showing it to anybody, put your retained striking cube into your jar and replace the lid. It is black if you struck full weight this round and white if you struck light, and it went under the cup before you touched anything else, so it cannot be anything but the truth.

Return the melt to the supply, refill the Silver bag to sixteen white and four black, redistribute one white and one black cube to each player, and pass the First Player marker clockwise.

\note{Why the jar, and not a reward you can see}{Any bonus paid out for having struck full weight would announce who struck full weight, and there goes the whole game. So honesty is not paid during the round at all. It is banked, in the dark, in a jar nobody may open, and settled at the end when secrecy no longer costs anything.

Your jar must hold exactly twelve cubes at the close, one for each round. A jar holding any other number has been tampered with.}

% ================= payouts =================
\sect{5}{Profit}

Score your action immediately in Phase 2, according to the current tier.

\begin{center}
\begin{tabular}{@{}lccc@{}}
\toprule
\headfont Trust & \headfont Full weight & \headfont False claim & \headfont Premium \\
\midrule
36 to 31 & 2 & 6 & 4 \\
30 to 21 & 2 & 5 & 3 \\
20 to 11 & 2 & 4 & 2 \\
10 to 1 & 2 & 3 & 1 \\
0 & \multicolumn{3}{c}{\itshape the coinage collapses; everybody loses} \\
\bottomrule
\end{tabular}
\end{center}

The premium for lying shrinks as the coinage degrades. Light coin is most profitable when the city can least afford it and barely worth the risk once the damage is done, which is the opposite of what everybody's instincts will tell them.

% ================= challenge odds =================
\sect{6}{Reading the Count}

You know the Light Count, and you know whether your own striking cube was black. Subtract it if so. The result is the \keyword{Remainder}, which contains the other four Houses' strikings and all five silver draws, including your own unknown one.

\begin{center}
\begin{tabular}{@{}cc@{}}
\toprule
\headfont Remainder & \headfont Chance any given House lied \\
\midrule
1 & 19\% \\
2 & 37\% \\
3 & 52\% \\
4 & 65\% \\
5 & 76\% \\
6 & 84\% \\
\bottomrule
\end{tabular}
\end{center}

Below a Remainder of 3, challenging is usually a poor bet. Above 4 it is close to a duty. Between the two, it depends entirely on what you think of the person you are looking at.

Remember that a bad shipment of silver puts a black cube in the melt exactly as deliberate short weight does. A House that was sold adulterated metal is indistinguishable from one that clipped its coin on purpose, and can never prove otherwise.

% ================= compact =================
\sect{7}{The Standard}

At any time between rounds, three or more players may sign the \keyword{Standard}: a public undertaking to strike full weight, checked against a shared set of scales.

\step{Dividend}{Members gain \keyword{+1 Profit} each round, paid whatever they struck. The Standard pays for belonging, not for behaving, because paying for behaviour would reveal it.}
\step{Tariff}{Non-members pay \keyword{1 Profit} to the bank each round. Trade on the city's terms costs more when you have not signed up to them.}
\step{Threshold}{The Standard only functions with three or more members. At two, all benefits cease at once.}

A member caught by challenge is sanctioned in steps, never all at once.

\step{First}{A Warning token.}
\step{Second}{Loss of the verification bonus for two rounds.}
\step{Third}{Expulsion, and the seat cannot be reclaimed.}

\note{Size decides everything}{With exactly three members, any single departure destroys the Standard for everybody, so every member holds a veto and knows it. With four, one may leave and it survives. The same act carries entirely different weight depending on how many are standing beside you.}

% ================= end =================
\sect{8}{Ending the Game}

The game ends after the twelfth round, or immediately if Trust reaches \keyword{0}.

If the coinage collapses, \keyword{every player loses}, including any who were eliminated earlier and any who were comfortably ahead.

Otherwise, each surviving player totals:

\bul{Profit accumulated.}
\bul{3 for every black cube in your jar.}
\bul{5 if still a member of a functioning Standard.}

Highest total wins. A player eliminated by Discredits scores nothing.

\note{On winning}{The winner is usually somebody who struck full weight most rounds, called an assay or two, and was believed. The second most common winner is somebody who lied twice, very early, and was never challenged for it. Both should feel earned.}

% ================= player aid =================
\clearpage
\thispagestyle{empty}
\begingroup
\setlength{\parskip}{0.3\baselineskip}
\linespread{1.0}\selectfont
\small

\begin{center}
  {\headfont\bfseries\large PLAYER AID\par}
\end{center}
{\color{rule}\hrule height 0.4pt}

{\headfont\bfseries The round\par}
\step{1.}{\keyword{Striking.} Fist into the melt, release one. Retained cube under the cup \emph{now}.}
\step{2.}{\keyword{Silver.} One cube from the Silver bag, \emph{unseen}, straight into the melt.}
\step{3.}{\keyword{Count.} Black cubes aloud.}
\step{4.}{\keyword{Assay.} One only, and only if you struck full weight.}
\step{5.}{\keyword{Erosion.} Trust drops by the Count. Retained cube into your \emph{jar}, lid on.}

{\headfont\bfseries Reading the count\par}
Remainder = Light Count minus your own striking cube, if it was black.

\begin{center}
\begin{tabular}{@{}lcccccc@{}}
\toprule
\headfont Remainder & 1 & 2 & 3 & 4 & 5 & 6 \\
\headfont Struck light & 19\% & 37\% & 52\% & 65\% & 76\% & 84\% \\
\bottomrule
\end{tabular}
\end{center}

Below 3, a poor bet. Above 4, close to a duty.

{\headfont\bfseries At the cup\par}
\begin{center}
\begin{tabular}{@{}ll@{}}
\toprule
\keyword{WHITE} under the cup & they played black. \keyword{They lied.} \\
\keyword{BLACK} under the cup & they struck full weight. \\
\bottomrule
\end{tabular}
\end{center}

{\headfont\bfseries Payouts\par}
\begin{center}
\begin{tabular}{@{}lcccc@{}}
\toprule
\headfont Trust & 36--31 & 30--21 & 20--11 & 10--1 \\
\headfont Honest & 2 & 2 & 2 & 2 \\
\headfont False & 6 & 5 & 4 & 3 \\
\bottomrule
\end{tabular}
\end{center}

{\color{rule}\hrule height 0.4pt}
\begin{center}
{\itshape Assay correctly: they lose the action and take a Discredit; you gain 3.\\
Assay wrongly: you take the Discredit. Three Discredits and you are ruined.\par}
\end{center}
\endgroup

\end{document}
